\chapter{Grammars and Pushdown Automata}\label{ch:grammars}

\begin{center}\itshape
To understand a sentence is to know how it was built;\\
a grammar is a recipe, and a parser cooks it backwards.\\
\upshape(the idea of this chapter, in our words)
\end{center}

More than two thousand years ago, the Sanskrit grammarian Pāṇini described the language of his time in a few thousand compact rules that rewrite symbols into other symbols, a system so precise that computer scientists still study it. In 1956 the young linguist Noam Chomsky, teaching at MIT, proposed to describe the sentences of English in the same spirit, by rules such as ``a sentence is a noun phrase followed by a verb phrase'', and classified such rule systems by their power. His paper \emph{Three models for the description of language} appeared in a journal on information theory, and it was read by engineers as much as by linguists.

Three years later, at the first international conference on information processing, held in Paris in June 1959 under the auspices of UNESCO, John Backus of IBM presented the syntax of a new programming language, the international algebraic language that became ALGOL, using what he called metalinguistic formulae. Peter Naur took up the notation as editor of the \emph{Report on the Algorithmic Language ALGOL~60}, signed by thirteen authors, which defines the whole syntax of the language in a few pages of rules like
\[
\langle\text{if statement}\rangle\ ::=\ \langle\text{if clause}\rangle\,\langle\text{unconditional statement}\rangle.
\]
The notation came to be called the Backus normal form. In 1964 Donald Knuth wrote to the \emph{Communications of the ACM} that it should rather be called the Backus--Naur form, since it is ``not a normal form'' in any usual sense; in 1967 Peter Ingerman suggested ``Pāṇini--Backus form'' in the same journal. Every programming language since has been defined by such a grammar, and every compiler begins by recovering, from a string of characters, the tree of rules that produced it.

The rules of ALGOL~60 had a surprising number of subtleties. Knuth's 1967 paper \emph{The remaining trouble spots in ALGOL~60} lists ambiguities and oddities in the report, and the rule displayed above is there for a reason we shall meet in Exercise~\ref{ex:dangling}: it keeps the program \texttt{if b then if c then x else y} from having two meanings.

\begin{think}
A single counter lets you check that a word has the form $a^nb^n$. A stack of plates, on which you may only put a plate on top or take the top one off, lets you check that brackets of several kinds are well nested, like \texttt{([()]())}. Can a stack check palindromes? Can it check $a^nb^nc^n$?
\end{think}

\section{Context-free grammars}\label{sec:cfg}

\begin{definition}\label{def:cfg}
A \emph{context-free grammar} is $G=(V,\Sigma,R,S)$, where $V$ is a finite set of \emph{variables}, $\Sigma$ a disjoint alphabet of \emph{terminals}, $S\in V$ the \emph{start variable}, and $R$ a finite set of \emph{rules} $A\to\alpha$ with $A\in V$ and $\alpha\in(V\cup\Sigma)^*$. If $A\to\alpha$ is a rule, then $\beta A\gamma\Rightarrow\beta\alpha\gamma$ for all strings $\beta,\gamma$; we write $\Rightarrow^*$ for a sequence of such steps, a \emph{derivation}. The \emph{language} of $G$ is $L(G)=\{w\in\Sigma^*:S\Rightarrow^*w\}$. A language is \emph{context-free} if some grammar generates it.
\end{definition}

``Context-free'' means that a variable is rewritten whatever surrounds it. For instance $S\to aSb\mid\varepsilon$ (two rules, separated by a bar) generates $\{a^nb^n\}$, which is not regular (Proposition~\ref{prop:nonregular}); $S\to aSa\mid bSb\mid a\mid b\mid\varepsilon$ generates the palindromes; and
\[
E\to E+T\mid T,\qquad T\to T*F\mid F,\qquad F\to(E)\mid a
\]
generates the arithmetic expressions built from $a$ with $+$, $*$ and parentheses.

A derivation can rewrite variables in many orders; what matters is the \emph{parse tree}: a rooted ordered tree whose internal nodes are labelled by variables, whose leaves are labelled by terminals or $\varepsilon$, and in which the children of a node labelled $A$ are labelled, from left to right, by the symbols of the right-hand side of a rule $A\to\alpha$ (a single leaf $\varepsilon$ if $\alpha=\varepsilon$). Its \emph{yield} is the word read on its leaves. A derivation is \emph{leftmost} if it always rewrites the leftmost variable.

\begin{proposition}\label{prop:parsetree}
For a variable $A$ and a word $w\in\Sigma^*$, the following are equivalent: $A\Rightarrow^*w$; there is a parse tree with root $A$ and yield $w$; there is a leftmost derivation $A\Rightarrow^*w$. Moreover, leftmost derivations of $w$ from $A$ and parse trees with root $A$ and yield $w$ are in one-to-one correspondence.
\end{proposition}

\begin{proof}
From a derivation we build a tree by induction on its length: the first step $A\Rightarrow X_1\cdots X_k$ uses a rule; the rest of the derivation rewrites each $X_i$ independently (the steps acting inside the part coming from $X_i$ form a derivation $X_i\Rightarrow^*w_i$, with $w=w_1\cdots w_k$), and by induction each gives a tree; hang them below a root $A$. From a tree with root $A$ and children $X_1,\dots,X_k$ we get a leftmost derivation: apply the rule at the root, then the leftmost derivation of the first subtree, then of the second, and so on; this is leftmost since everything to the left is already terminal. These two constructions are inverse to each other between trees and leftmost derivations: a leftmost derivation is determined by the sequence of rules it applies, and this sequence is the preorder (root, then subtrees from left to right) of the rules used in the tree.
\end{proof}

\begin{figure}[ht]
\centering
\fig{ch02-parsetree}
\caption{The parse tree of $a+a*a$ in the expression grammar. The product $a*a$ sits deeper in the tree than the sum, so it is computed first: the grammar encodes the priority of $*$ over $+$.}
\label{fig:parsetree}
\end{figure}

\begin{proposition}\label{prop:dyck}
The grammar $S\to(S)S\mid\varepsilon$ generates exactly the \emph{Dyck words}: the words over $\{(,)\}$ with as many opening as closing brackets, in which no prefix has more closing than opening brackets.
\end{proposition}

\begin{proof}
Let $h(w)$ be the number of opening minus closing brackets in $w$; Dyck words are those with $h(w)=0$ and $h(u)\ge0$ for every prefix $u$. By induction on the derivation, every generated word is a Dyck word: $\varepsilon$ is, and if $x,y$ are, the prefixes of $(x)y$ have heights $1+h(u)\ge1$ for prefixes $(u$ with $u$ a prefix of $x$, then $0$ after $(x)$, then $h(u)\ge0$ inside $y$. Conversely, let $w\ne\varepsilon$ be a Dyck word. It starts with $($, since its first prefix has height $\ge0$. Let $(x)$ be its shortest nonempty prefix of height $0$, so $w=(x)y$. The heights of the prefixes of $x$ are those of the corresponding prefixes of $w$ minus $1$, which are $\ge0$ by minimality, and $h(x)=0$; so $x$ is a Dyck word, and so is $y$ (its prefixes have the heights of the prefixes of $w$ that contain $(x)$). By induction on length, $S\Rightarrow^*x$ and $S\Rightarrow^*y$, and $S\Rightarrow(S)S\Rightarrow^*(x)y$.
\end{proof}

The decomposition $w=(x)y$ in the proof is unique, which will show that the grammar is unambiguous (Exercise~\ref{ex:dyckcount}). Every regular language is context-free: given a DFA, take a variable $A_q$ for each state, rules $A_q\to aA_{\delta(q,a)}$, and $A_q\to\varepsilon$ for $q\in F$; then $A_{q}\Rightarrow^*w$ exactly when $\delta^*(q,w)\in F$, by induction on $|w|$.

\section{Ambiguity}\label{sec:ambiguity}

\begin{definition}\label{def:ambiguous}
A grammar is \emph{ambiguous} if some word has two different parse trees, equivalently (Proposition~\ref{prop:parsetree}) two different leftmost derivations.
\end{definition}

The grammar $E\to E+E\mid E*E\mid a$ generates the same expressions without parentheses as the grammar above, but $a+a*a$ has two parse trees: one computes $(a+a)*a$, the other $a+(a*a)$. For a compiler the tree \emph{is} the meaning, so ambiguity is a bug. It can often be removed, by layering the grammar as in Figure~\ref{fig:parsetree}, but not always: Rohit Parikh showed in 1961 that $\{a^ib^jc^k:i=j\text{ or }j=k\}$ is \emph{inherently ambiguous}, generated only by ambiguous grammars, a result we quote. Whether a given grammar is ambiguous cannot be decided by any algorithm (Chapter~4).

\section{Chomsky normal form and parsing}\label{sec:cnf}

\begin{definition}\label{def:cnf}
A grammar is in \emph{Chomsky normal form} if every rule has the form $A\to BC$ with $B,C$ variables or $A\to a$ with $a$ a terminal.
\end{definition}

Such a grammar cannot generate $\varepsilon$, and in it a parse tree for a word of length $n\ge1$ has exactly $n$ rules of the form $A\to a$ and $n-1$ of the form $A\to BC$, since each binary rule increases the number of variables by one.

\begin{theorem}[Chomsky, 1959]\label{thm:cnf}
For every context-free grammar $G$ there is a grammar in Chomsky normal form generating $L(G)\setminus\{\varepsilon\}$.
\end{theorem}

\begin{proof}
We transform $G$ in four steps, each preserving the set of nonempty words derived from every variable.

\emph{Removing $\varepsilon$-rules.} A variable is \emph{nullable} if it derives $\varepsilon$; the nullable variables are computed by starting with those having a rule $A\to\varepsilon$ and adding $A$ whenever it has a rule $A\to X_1\cdots X_k$ with all $X_i$ nullable, until nothing changes (a derivation of $\varepsilon$ has a parse tree all of whose nodes are nullable, by induction on its height). For every rule $A\to X_1\cdots X_k$, add all the rules obtained by deleting some of the nullable $X_i$, except the empty right-hand side; then delete all rules $A\to\varepsilon$. The new rules are simulated by old ones followed by derivations of $\varepsilon$, so no new word appears. Conversely, take a parse tree in the old grammar with root $A$ and nonempty yield; at its root $A\to X_1\cdots X_k$, the children $X_i$ whose subtrees yield $\varepsilon$ are nullable, and dropping them uses an allowed new rule (at least one child remains, since the yield is nonempty); by induction on the size, the remaining subtrees, whose yields are nonempty, can be rebuilt in the new grammar.

\emph{Removing unit rules} $A\to B$. Compute, by reachability in the graph of unit rules, the pairs with $A\Rightarrow^*B$ through unit rules only (including $A=B$). For each such pair and each non-unit rule $B\to\alpha$, add $A\to\alpha$; then delete the unit rules. A parse tree in the old grammar is turned into one in the new grammar by collapsing every maximal chain of nodes linked by unit rules into its top node, and conversely the new rules are simulated by chains of unit rules.

\emph{Shortening rules.} Replace each rule $A\to X_1X_2\cdots X_k$ with $k\ge3$ by $A\to X_1A_1$, $A_1\to X_2A_2$, \dots, $A_{k-2}\to X_{k-1}X_k$, with new variables $A_i$ used nowhere else; a use of the old rule corresponds to a use of the chain.

\emph{Separating terminals.} In each rule $A\to X_1X_2$, replace a terminal $X_i=a$ by a new variable $T_a$ with the single rule $T_a\to a$.

The remaining rules are of the forms $A\to BC$ and $A\to a$, and the variable $S$ derives exactly the nonempty words of $L(G)$.
\end{proof}

The normal form makes parsing a dynamic programme, discovered independently by John Cocke, Tadao Kasami and Daniel Younger in the 1960s.

\begin{theorem}[Cocke--Kasami--Younger]\label{thm:cyk}
Given a grammar in Chomsky normal form with $r$ rules and a word $w=a_1\cdots a_n$, one can decide whether $w\in L(G)$, and find a parse tree, in time $O(rn^3)$.
\end{theorem}

\begin{proof}
For $1\le i\le j\le n$ let $T_{ij}$ be the set of variables $A$ with $A\Rightarrow^*a_i\cdots a_j$. A parse tree of a single letter is a rule $A\to a_i$, so $T_{ii}=\{A:(A\to a_i)\in R\}$. For $i<j$, the root rule of a parse tree is $A\to BC$, and $B$ yields a nonempty prefix $a_i\cdots a_k$ and $C$ the nonempty rest, so
\[
T_{ij}=\{A:\ (A\to BC)\in R,\ B\in T_{ik},\ C\in T_{k+1,j}\text{ for some }i\le k<j\}.
\]
Computing the tables by increasing $j-i$ takes $O(n)$ choices of $k$ and $O(r)$ rules for each of the $O(n^2)$ cells. Then $w\in L(G)$ if and only if $S\in T_{1n}$, and remembering for each $A\in T_{ij}$ one successful pair $(k,A\to BC)$ lets one rebuild a parse tree.
\end{proof}

\begin{figure}[ht]
\centering
\fig{ch02-cyk}
\caption{The table of Theorem~\ref{thm:cyk} for the word $aaabbb$ and the grammar $S\to AB\mid AX$, $X\to SB$, $A\to a$, $B\to b$, a Chomsky normal form of $S\to aSb\mid ab$. The cell in row $\ell$ above a factor of length $\ell$ lists the variables deriving that factor; $S$ in the top cell means that the word is accepted.}
\label{fig:cyk}
\end{figure}

Leslie Valiant showed in 1975 that the cubic time can be lowered to the time of Boolean matrix multiplication, $O(n^{2.38})$ with the current exponents (Tome~C2, Chapter~9). Compilers do much better, in linear time, by restricting the grammars: the $LR(k)$ grammars of Knuth (1965) are parsed deterministically, reading the program once from left to right with $k$ letters of lookahead.

\section{Pushdown automata}\label{sec:pda}

A finite automaton with a stack is the machine that matches context-free grammars, as finite automata match regular expressions.

\begin{definition}\label{def:pda}
A \emph{pushdown automaton} (PDA) is $P=(Q,\Sigma,\Gamma,\delta,q_0,F)$, with a finite \emph{stack alphabet} $\Gamma$ and a transition function giving, for each state $q$, each $x\in\Sigma\cup\{\varepsilon\}$ and each $t\in\Gamma\cup\{\varepsilon\}$, a finite set $\delta(q,x,t)$ of pairs $(q',t')$ with $t'\in\Gamma\cup\{\varepsilon\}$. Such a pair is a move: in state $q$, read $x$ from the input (nothing if $x=\varepsilon$), pop $t$ from the top of the stack (nothing if $t=\varepsilon$), push $t'$ (nothing if $t'=\varepsilon$), and go to state $q'$. The PDA \emph{accepts} $w$ if some sequence of moves, starting in $q_0$ with an empty stack, reads all of $w$ and ends in a state of $F$.
\end{definition}

For instance, a PDA for the palindromes pushes the letters it reads; at some point it guesses that it has reached the middle (possibly skipping one middle letter), and from then on it pops a letter for each letter read, checking that they are equal; it accepts if the stack is empty at the end (detected by a special symbol $\$$ pushed at the start). The guess is essential: the deterministic PDAs, which never have a choice, recognise a strictly smaller class that does not contain the palindromes, a theorem we quote.

\begin{theorem}\label{thm:cfgtopda}
Every context-free language is recognised by a PDA.
\end{theorem}

\begin{proof}
Let $G=(V,\Sigma,R,S)$. The PDA simulates leftmost derivations on its stack. We allow ourselves to push a whole string $X_1\cdots X_k$ in one move; this is implemented by $k$ moves through new intermediate states, pushing $X_k$ first so that $X_1$ ends on top. The stack alphabet is $V\cup\Sigma\cup\{\$\}$, and the states are $q_0$, $q$ and an accepting state $f$. The first move pushes $S\$$ and goes to $q$. In state $q$:
\begin{itemize}
\item if the top is a variable $A$, for each rule $A\to\alpha$ there is a move that pops $A$ and pushes $\alpha$, reading nothing;
\item if the top is a terminal $a$, there is a move that reads $a$ and pops it;
\item if the top is $\$$, there is a move that pops it and goes to $f$.
\end{itemize}
\emph{Invariant.} If, in state $q$, the PDA has read $u$ and its stack is $\beta\$$ ($\beta$ read from the top), then $S\Rightarrow^*u\beta$ by a leftmost derivation. This holds after the first move ($u=\varepsilon$, $\beta=S$) and is preserved: expanding the top variable $A$ of $\beta=A\gamma$ performs the derivation step $uA\gamma\Rightarrow u\alpha\gamma$, which is leftmost since $u$ is terminal, and matching a terminal changes $u$ to $ua$ and $\beta=a\gamma$ to $\gamma$, leaving $u\beta$ unchanged. So if the PDA reaches $f$ after reading $w$, then $\beta=\varepsilon$ when it popped $\$$, and $S\Rightarrow^*w$.

\emph{Conversely}, let $S=\gamma_0\Rightarrow\gamma_1\Rightarrow\cdots\Rightarrow\gamma_m=w$ be a leftmost derivation, and write $\gamma_i=u_iA_i\beta_i$ with $u_i$ terminal. We show by induction on $i$ that the PDA can reach state $q$ having read $u_i$ with stack $A_i\beta_i\$$: from there, expanding $A_i$ by the rule used in step $i+1$ gives the stack $\alpha\beta_i\$$, and matching the terminal prefix of $\alpha\beta_i$ against the input (which is the next part of $w$, since $w$ extends $u_{i+1}$) reaches the configuration for $i+1$. After the last step, the stack holds the terminal rest of $w$ followed by $\$$; matching it and popping $\$$ accepts $w$.
\end{proof}

\begin{theorem}\label{thm:pdatocfg}
Every language recognised by a PDA is context-free.
\end{theorem}

\begin{proof}
First we normalise the PDA $P$, without changing its language: (i) it has a single accepting state $f$; (ii) it empties its stack before accepting; (iii) each move either pushes one symbol or pops one symbol, but not both and not neither. For (i) and (ii), push a new bottom symbol $\$$ at the start, add moves from the old accepting states to a new state that pops any symbol, and from there pop $\$$ into the single accepting state $f$. For (iii), split a move that pops and pushes into a pop followed by a push through a new state, and replace a move that does neither by a push of a new dummy symbol followed by its pop.

For each pair of states $p,q$ we introduce a variable $A_{pq}$, intended to generate the words $x$ that can take $P$ from state $p$ with empty stack to state $q$ with empty stack. Such a computation never looks below the stack it started with, so the same words take $P$ from $p$ with any stack $\gamma$ to $q$ with the same stack $\gamma$, never touching $\gamma$. The rules are:
\begin{itemize}
\item $A_{pp}\to\varepsilon$ for each $p$;
\item $A_{pq}\to A_{pr}A_{rq}$ for all $p,q,r$;
\item $A_{pq}\to xA_{rs}y$ whenever $(r,t)\in\delta(p,x,\varepsilon)$ and $(q,\varepsilon)\in\delta(s,y,t)$ for some stack symbol $t$ and some $x,y\in\Sigma\cup\{\varepsilon\}$.
\end{itemize}
The start variable is $A_{q_0f}$.

\emph{If $A_{pq}\Rightarrow^*x$, then $x$ takes $P$ from $p$ with empty stack to $q$ with empty stack.} By induction on the length of the derivation: the first rule is $A_{pp}\to\varepsilon$ (nothing to do), or $A_{pq}\to A_{pr}A_{rq}$ (run the two computations one after the other), or $A_{pq}\to xA_{rs}y$, with $x'$ derived from $A_{rs}$: read $x$ pushing $t$, run the computation for $x'$ from $r$ to $s$ above $t$, then read $y$ popping $t$.

\emph{If $x$ takes $P$ from $p$ with empty stack to $q$ with empty stack, then $A_{pq}\Rightarrow^*x$.} By induction on the number of moves. With no move, $p=q$, $x=\varepsilon$ and $A_{pp}\to\varepsilon$. Otherwise the first move pushes (the stack is empty) and the last move pops (the stack ends empty). If the stack becomes empty at some intermediate moment, in some state $r$, the computation splits into two shorter ones, for words $x_1x_2=x$, and $A_{pq}\to A_{pr}A_{rq}\Rightarrow^*x_1x_2$. If not, the symbol $t$ pushed by the first move stays at the bottom until the last move, which pops it: the first move goes from $p$ to some $r$ reading some $x_0$, the last from some $s$ to $q$ reading some $y_0$, and in between a shorter computation takes $r$ to $s$ above $t$, reading $x'$ with $x=x_0x'y_0$. By induction $A_{rs}\Rightarrow^*x'$, and the third kind of rule gives $A_{pq}\Rightarrow x_0A_{rs}y_0\Rightarrow^*x$.

Since $P$ accepts $w$ exactly when $w$ takes it from $q_0$ with empty stack to $f$ with empty stack, the grammar generates $L(P)$.
\end{proof}

\section{What one stack cannot do}\label{sec:cflpumping}

\begin{lemma}[Bar-Hillel, Perles, Shamir, 1961]\label{lem:cflpumping}
For every context-free language $L$ there is $p$ such that every $w\in L$ with $|w|\ge p$ can be written $w=uvxyz$ with $|vy|\ge1$, $|vxy|\le p$, and $uv^ixy^iz\in L$ for all $i\ge0$.
\end{lemma}

\begin{proof}
Take a grammar in Chomsky normal form for $L\setminus\{\varepsilon\}$ (Theorem~\ref{thm:cnf}) with $m$ variables, and let $p=2^m$. In a parse tree, call the \emph{height} the largest number of variables on a path from the root to a leaf. A parse tree of height $h$ yields at most $2^{h-1}$ letters, by induction ($A\to a$ has height $1$ and yields one letter; $A\to BC$ adds one level above two subtrees). So a parse tree of $w$ with $|w|\ge2^m$ has height at least $m+1$.

Take a longest path from the root, and look at its last $m+1$ variables: two of them carry the same variable $R$. Let the upper one be the root of a subtree yielding $vxy$ and the lower one the root of a subtree yielding $x$, so that $w=uvxyz$. Replacing the lower subtree by a copy of the upper one, $i-1$ times, gives parse trees of $uv^ixy^iz$; replacing the upper subtree by the lower one gives $uxz$. The upper $R$ uses a rule $R\to BC$, and the lower $R$ lies in the subtree of $B$ or of $C$; the other one yields a nonempty word (there are no $\varepsilon$-rules), which is part of $v$ or of $y$, so $|vy|\ge1$. Finally, every path down from the upper $R$ has at most $m+1$ variables, since the chosen path is a longest one and its part below the upper $R$ has at most $m+1$; so $|vxy|\le2^m=p$.
\end{proof}

\begin{proposition}\label{prop:noncf}
The languages $\{a^nb^nc^n:n\ge0\}$ and $\{ww:w\in\{a,b\}^*\}$ are not context-free.
\end{proposition}

\begin{proof}
Pump $a^pb^pc^p$: since $|vxy|\le p$, the factor $vxy$ meets at most two of the three blocks, so $uv^2xy^2z$ has more letters of one or two kinds and the same number of the third (and if $v$ or $y$ straddles two blocks, the letters are not even in order). For $\{ww\}$, pump $s=a^pb^pa^pb^p$ down, to $uxz$. A word $ww$ with $w\ne\varepsilon$ has either one maximal block of equal letters or at least four: if $w$ starts and ends with the same letter, $ww$ has $2k-1$ blocks when $w$ has $k$, and $k\ge2$ unless $w$ is a power of one letter; otherwise $ww$ has $2k\ge4$ blocks. The window $vxy$, of length at most $p$, meets at most two adjacent blocks of $s$ and empties at most one of them; so $uxz$ has two, three or four blocks. With two or three it is not of the form $ww$. With four it is $a^ib^ja^kb^l$ with $(i,j)\ne(k,l)$: if the window meets only the first two blocks, $i+j<2p=k+l$; if only the last two, $k+l<2p=i+j$; if the middle two, $j<p=l$ or $k<p=i$. A word $a^ib^ja^kb^l$ with four nonempty blocks is of the form $ww$ only if $w=a^ib^j=a^kb^l$. So $uxz\notin\{ww\}$.
\end{proof}

\begin{proposition}\label{prop:cflclosure}
Context-free languages are closed under union, concatenation, star, homomorphisms, and intersection with a regular language. They are not closed under intersection or complement.
\end{proposition}

\begin{proof}
For grammars with start variables $S_1,S_2$ and disjoint variables, $S\to S_1\mid S_2$, $S\to S_1S_2$ and $S\to SS_1\mid\varepsilon$ give the union, the concatenation and the star; replacing each terminal $a$ by the word $h(a)$ in all rules gives $h(L)$. For a regular language, run a DFA alongside a PDA: the states are pairs, a move reading $x$ updates the DFA component by $x$ (or leaves it if $x=\varepsilon$), and a pair accepts when both components do; a run of the product is a run of each machine on the same word. Finally $\{a^nb^nc^m\}$ and $\{a^mb^nc^n\}$ are context-free, while their intersection $\{a^nb^nc^n\}$ is not; since $L\cap M$ is the complement of the union of the complements, the context-free languages are not closed under complement either.
\end{proof}

\begin{remark}
Over a one-letter alphabet, context-free languages are regular. This follows from Parikh's theorem (1961, published 1966), which says that the vectors counting the letters of the words of a context-free language form a finite union of sets $\{v_0+n_1v_1+\cdots+n_kv_k:n_i\ge0\}$; we quote it. So $\{a^{n^2}\}$ and $\{a^p:p\text{ prime}\}$ are not context-free.
\end{remark}

\begin{proposition}\label{prop:cfgdecide}
Given a grammar $G$, one can decide whether $L(G)$ is empty and whether it is finite.
\end{proposition}

\begin{proof}
Call a variable \emph{generating} if it derives some terminal word; the generating variables are found by the same saturation as the nullable ones in Theorem~\ref{thm:cnf}, starting from the rules whose right-hand sides are terminal, and $L(G)\ne\emptyset$ if and only if $S$ is generating. For finiteness, see Exercise~\ref{ex:finite}.
\end{proof}

In contrast, deciding whether two grammars generate the same language, or whether a grammar generates all words, is impossible for any algorithm; Chapter~4 proves it.

\begin{beauty}
A context-free grammar is a finite description of an infinite set of trees, and the theorems of this chapter are all statements about trees: the normal form makes them binary, the cubic parser fills a triangle with the subtrees of every factor, the pumping lemma finds a variable repeated on a long branch and grafts the branch onto itself. The pushdown automaton is the same object seen from the inside: a stack is exactly what one needs to walk a tree without losing one's way, and the grammar of Theorem~\ref{thm:pdatocfg} rebuilds the tree from the moments when the stack returns to its earlier height. That Pāṇini's rules, Chomsky's linguistics and Backus's compilers turned out to be one mathematical object is the kind of coincidence that is not one.
\end{beauty}

\section{Exercises}

\begin{exercise}[$\star\star$]\label{ex:dyckcount}
Show that the grammar $S\to(S)S\mid\varepsilon$ is unambiguous, and deduce that the number $D_n$ of Dyck words of length $2n$ satisfies $D_0=1$ and $D_{n+1}=\sum_{k=0}^nD_kD_{n-k}$, the recurrence of the Catalan numbers.
\end{exercise}

\begin{exercise}[$\star\star$]\label{ex:equalab}
Show that $S\to aSbS\mid bSaS\mid\varepsilon$ generates the words over $\{a,b\}$ with as many $a$'s as $b$'s. Is this grammar ambiguous?
\end{exercise}

\begin{exercise}[$\star$]\label{ex:catalanparse}
Show that in the grammar $E\to E+E\mid a$, the word $a+a+\cdots+a$ with $n+1$ letters $a$ has $D_n$ parse trees, with $D_n$ as in Exercise~\ref{ex:dyckcount}.
\end{exercise}

\begin{exercise}[$\star\star$]\label{ex:dangling}
Show that the grammar $S\to\texttt{i}S\mid\texttt{i}S\texttt{e}S\mid\texttt x$ (with \texttt i for ``\texttt{if b then}'', \texttt e for ``\texttt{else}'' and \texttt x for any other statement) is ambiguous. Then show that the grammar of ALGOL~60, in this notation
\[
S\to U\mid\texttt iU\mid\texttt iU\texttt eS,\qquad U\to\texttt x\mid\texttt bS\texttt d
\]
(with \texttt b and \texttt d for \texttt{begin} and \texttt{end}), is unambiguous.
\end{exercise}

\begin{exercise}[$\star$]\label{ex:cykcount}
Modify the algorithm of Theorem~\ref{thm:cyk} to count the parse trees of $w$ in time $O(rn^3)$ (counting arithmetic operations). Deduce an algorithm deciding whether a given word has two parse trees.
\end{exercise}

\begin{exercise}[$\star$]\label{ex:noncfmore}
Show that $\{a^ib^jc^k:i\le j\le k\}$ and $\{a^{n^2}:n\ge0\}$ are not context-free, using Lemma~\ref{lem:cflpumping}.
\end{exercise}

\begin{exercise}[$\star\star\star$]\label{ex:complementww}
Show that the complement of $\{ww:w\in\{a,b\}^*\}$ is context-free, although $\{ww\}$ is not.
\end{exercise}

\begin{exercise}[$\star\star$]\label{ex:finite}
Let $G$ be in Chomsky normal form with all its variables generating and reachable from $S$. Draw an arrow $A\to B$ and $A\to C$ for every rule $A\to BC$. Show that $L(G)$ is infinite if and only if this graph has a cycle.
\end{exercise}

\begin{exercise}[$\star$]\label{ex:pdaanbn}
Describe a PDA, with the conventions of Definition~\ref{def:pda}, for the words over $\{a,b\}$ with as many $a$'s as $b$'s.
\end{exercise}

\begin{exercise}[$\star\star$]\label{ex:linearlength}
Show that if $G$ is in Chomsky normal form with $m$ variables and $L(G)$ is nonempty, then it contains a word of length at most $2^{m-1}$, and that this bound is attained.
\end{exercise}

\section{Hints}

\begin{hint}{ex:dyckcount}
In the proof of Proposition~\ref{prop:dyck}, the first rule of a parse tree of $w\ne\varepsilon$ is forced, and so is the split $w=(x)y$.
\end{hint}
\begin{hint}{ex:equalab}
For the converse, look at the first return to height $0$ of the function counting $a$'s minus $b$'s. For ambiguity, try $abab$.
\end{hint}
\begin{hint}{ex:catalanparse}
The root rule $E\to E+E$ chooses which $+$ is computed last.
\end{hint}
\begin{hint}{ex:dangling}
For the first grammar, try $\texttt{iixex}$. For the second, show by induction on length that the first symbols of a word determine which rule was used at the root and where each part ends.
\end{hint}
\begin{hint}{ex:cykcount}
Replace sets by counts $N_{ij}(A)$ and the union by a sum of products.
\end{hint}
\begin{hint}{ex:noncfmore}
For the first, pump $a^pb^pc^p$ up or down according to which blocks $v$ and $y$ meet. For the second, compare $p^2+|vy|$ with $(p+1)^2$.
\end{hint}
\begin{hint}{ex:complementww}
A word of even length is not of the form $ww$ when its two halves differ at some position; regroup the word around the two different letters.
\end{hint}
\begin{hint}{ex:finite}
A cycle through $A$ gives $A\Rightarrow^*xAy$ with $xy\ne\varepsilon$. Without cycles, the height of parse trees is bounded.
\end{hint}
\begin{hint}{ex:pdaanbn}
Keep on the stack the surplus of the letter that is ahead.
\end{hint}
\begin{hint}{ex:linearlength}
A parse tree of minimal height repeats no variable on a branch. For the bound, use the rules $A_1\to A_2A_2$, \dots, $A_{m-1}\to A_mA_m$, $A_m\to a$.
\end{hint}

\section{Solutions}

\begin{solution}{ex:dyckcount}
\idea Unambiguity means that the parse tree can be reconstructed from the word, and the proof of Proposition~\ref{prop:dyck} already reconstructs it: the split $(x)y$ is forced.
\writeup We show by induction on $|w|$ that every Dyck word has exactly one parse tree. For $w=\varepsilon$, a parse tree must use $S\to\varepsilon$ at the root (the other rule produces at least two letters). For $w\ne\varepsilon$ the root rule must be $S\to(S)S$, and the subtrees yield $x,y$ with $w=(x)y$, where $x$ and $y$ are Dyck words (Proposition~\ref{prop:dyck} applied to the subtrees). Then $(x)$ is a prefix of height $0$, and it is the shortest nonempty one: a shorter nonempty prefix $(u$ of height $0$ would give a prefix $u$ of $x$ with height $-1$. So $x$ and $y$ are determined by $w$, and by induction their subtrees are unique.
Counting by the unique decomposition, a Dyck word of length $2(n+1)$ is $(x)y$ with $|x|=2k$ and $|y|=2(n-k)$ for a unique $k\in\{0,\dots,n\}$, and every such pair gives one; hence $D_{n+1}=\sum_kD_kD_{n-k}$. These are the Catalan numbers, $D_n=\frac1{n+1}\binom{2n}n$.
\end{solution}

\begin{solution}{ex:equalab}
\idea Track the difference between the numbers of $a$'s and $b$'s along the word. A word with difference $0$ that starts with $a$ first returns to $0$ at a $b$, which splits it as $axby$.
\writeup Every generated word has as many $a$'s as $b$'s, by induction on derivations. Conversely, let $w\ne\varepsilon$ be balanced and $d(u)$ the number of $a$'s minus the number of $b$'s of a prefix $u$. Say $w$ starts with $a$ (the other case is symmetric). Let $u$ be the shortest nonempty prefix with $d(u)=0$; it ends with $b$ (just before, $d=1$), so $u=axb$, and $x$ is balanced; and $y$, with $w=uy$, is balanced. By induction on length, $S\Rightarrow^*x$ and $S\Rightarrow^*y$, so $S\Rightarrow aSbS\Rightarrow^*w$.
The grammar is ambiguous: $abab$ is obtained with $x=\varepsilon$, $y=ab$ (the split above), but also as $a\,S\,b\,S$ with the first $S$ deriving $ba$ and the second $\varepsilon$, since $S\Rightarrow bSaS\Rightarrow ba$.
\end{solution}

\begin{solution}{ex:catalanparse}
\idea The root of a parse tree splits the sum at one of its $n$ plus signs; the two sides are again sums.
\writeup Let $P_n$ be the number of parse trees of $a(+a)^n$. For $n=0$ the only tree is $E\to a$, so $P_0=1$. For $n\ge1$ the root uses $E\to E+E$, the two subtrees yield $a(+a)^k$ and $a(+a)^{n-1-k}$ for some $k\in\{0,\dots,n-1\}$ (the $+$ of the root rule being the $(k+1)$st one), and the subtrees are arbitrary. So $P_n=\sum_{k=0}^{n-1}P_kP_{n-1-k}$, the recurrence of Exercise~\ref{ex:dyckcount}, with the same initial value; hence $P_n=D_n$.
\end{solution}

\begin{solution}{ex:dangling}
\idea In the first grammar, an \texttt e after two \texttt i's can belong to either. ALGOL~60 forbids an \texttt i directly after an \texttt i: the statement after \texttt{then} must be unconditional, so an \texttt{if} inside it must be wrapped in \texttt{begin}\dots\texttt{end}, and the \texttt{end} tells where it stops.
\writeup The word $\texttt{iixex}$ has two parse trees: $S\to\texttt iS$ with $S\Rightarrow\texttt{ixex}$ by $S\to\texttt iS\texttt eS$ (the \texttt e belongs to the inner \texttt i), and $S\to\texttt iS\texttt eS$ with the first $S\Rightarrow\texttt{ix}$ and the second $S\Rightarrow\texttt x$ (the \texttt e belongs to the outer one).
For the ALGOL grammar we show, by induction on the length of a word $w$, that $w$ has at most one parse tree from $S$ and at most one from $U$. From $U$: if $w$ starts with \texttt x, the rule is $U\to\texttt x$ and $w=\texttt x$; if it starts with \texttt b, the rule is $U\to\texttt bS\texttt d$, $w=\texttt bw'\texttt d$, and $w'$ has at most one tree by induction. Note also that no word derived from $U$ is a proper prefix of another one: words from $U$ are \texttt x or have balanced \texttt b/\texttt d with the final \texttt d closing the first \texttt b, so the end of a word from $U$ is determined by reading it from the left (the first moment the count of \texttt b minus \texttt d returns to $0$). From $S$: if $w$ does not start with \texttt i, the rule is $S\to U$. If it starts with \texttt i, it is $\texttt iu$ or $\texttt iu\texttt ev$ with $u$ derived from $U$; the word $u$ is the unique prefix of the rest derived from $U$, so it is determined, and then either nothing follows (rule $S\to\texttt iU$) or \texttt e follows (rule $S\to\texttt iU\texttt eS$, with $v$ determined). In each case the rule and the parts are determined, and by induction so are the subtrees.
\end{solution}

\begin{solution}{ex:cykcount}
\idea Parse trees of $a_i\cdots a_j$ with root $A$ are classified by the root rule and the split point, and for each choice they are pairs of trees.
\writeup Let $N_{ij}(A)$ be the number of parse trees with root $A$ and yield $a_i\cdots a_j$. Then $N_{ii}(A)=1$ if $A\to a_i$ is a rule and $0$ otherwise, and $N_{ij}(A)=\sum_{(A\to BC)\in R}\sum_{k=i}^{j-1}N_{ik}(B)N_{k+1,j}(C)$, since a tree is determined by its root rule, the length of the yield of the left subtree, and the two subtrees. The table is filled in $O(rn^3)$ operations, and $w$ has $N_{1n}(S)$ parse trees; it is ambiguous for this grammar when this number is at least $2$. (The numbers can be exponentially large, with $O(n)$ digits; to decide only whether the count is at least $2$, cap all counts at $2$.)
\end{solution}

\begin{solution}{ex:noncfmore}
\idea Pumping adds letters to at most two adjacent blocks; choose the direction of pumping that breaks an inequality. For squares, pumping changes the length by at most $p$, which is too little to reach the next square.
\writeup Let $w=a^pb^pc^p$. The window $vxy$ meets at most two adjacent blocks. If $v$ or $y$ contains two kinds of letters, $uv^2xy^2z$ is not in $a^*b^*c^*$. Otherwise: if $vy$ contains an $a$, then pumping up ($i=2$) gives more $a$'s than $c$'s, or more $a$'s than $b$'s if no $b$ is added, since no $c$ is added; if $vy$ contains no $a$ but a $b$, pumping up gives more $b$'s than $c$'s unless a $c$ is also added, in which case pumping down ($i=0$) gives fewer $b$'s than $a$'s; if $vy$ contains only $c$'s, pumping down gives fewer $c$'s than $b$'s. In all cases the word leaves the language.
For $\{a^{n^2}\}$, take $w=a^{p^2}$; then $uv^2xy^2z$ has length $p^2+|vy|$ with $1\le|vy|\le p<2p+1$, strictly between $p^2$ and $(p+1)^2$, so it is not a square.
\end{solution}

\begin{solution}{ex:complementww}
\idea Every word of odd length is in the complement. A word of even length $2n$ is in it when its two halves differ at some position $i$: it reads $u\,a\,v\,u'\,b\,v'$ with $|u|=|u'|=i-1$ and $|v|=|v'|=n-i$ (or with $a,b$ exchanged). The constraint links lengths that are far apart, but the middle part $vu'$ can be cut differently, into a part as long as $u$ and a part as long as $v'$: then the word becomes a word of odd length with $a$ in its middle followed by a word of odd length with $b$ in its middle, and both are context-free.
\writeup Let $A\to aAa\mid aAb\mid bAa\mid bAb\mid a$ and $B$ the same with last rule $B\to b$, so that $A$ generates the words of odd length with middle letter $a$, and $B$ those with middle letter $b$. Let $S\to AB\mid BA\mid O$ with $O\to aE\mid bE$ and $E\to aaE\mid abE\mid baE\mid bbE\mid\varepsilon$ generating the words of odd length. We show $L(S)$ is the complement $K$ of $\{ww\}$.
If $s\in L(A)L(B)$, then $s=u\,a\,u''\,t''\,b\,t$ with $|u''|=|u|$ and $|t''|=|t|$. Its length is $2n$ with $n=|u|+|t|+1$, and its letter at position $|u|+1$ is $a$. The first factor has length $2|u|+1$, so position $n+|u|+1=(2|u|+1)+(|t|+1)$ of $s$ is position $|t|+1$ of the second factor $t''bt$, its middle letter $b$. So the two halves of $s$ differ at position $|u|+1$, and $s\in K$. The same holds for $L(B)L(A)$, and $L(O)\subseteq K$.
Conversely, let $s\in K$. If $|s|$ is odd, $s\in L(O)$. If $|s|=2n$ and the halves differ at position $i$, say $s_i=a$ and $s_{n+i}=b$, let $u=s_1\cdots s_{i-1}$, and split $s=s_1\cdots s_{2i-1}\cdot s_{2i}\cdots s_{2n}$. The first factor has length $2i-1$ and middle letter $s_i=a$, so it is in $L(A)$. The second has length $2n-2i+1$, and its middle position is $(2n-2i+2)/2=n-i+1$, which is position $2i-1+n-i+1=n+i$ of $s$, the letter $b$; so it is in $L(B)$. Similarly with $a,b$ exchanged. So $s\in L(S)$.
\end{solution}

\begin{solution}{ex:finite}
\idea A cycle in the graph is a variable that can reproduce itself with something on the sides, which is pumping. Without cycles, all branches of parse trees are short.
\writeup If the graph has a cycle through $A$, then $A\Rightarrow^*xAy$ for some $x,y$ with $xy\ne\varepsilon$: each arrow $A\to B$ along the cycle comes from a rule $A\to BC$ or $A\to CB$, and the other variable $C$, being generating, derives a nonempty terminal word (no $\varepsilon$ in Chomsky normal form). As $A$ is reachable and generating, $S\Rightarrow^*uAz$ and $A\Rightarrow^*v$ for terminal words, so $S\Rightarrow^*ux^kvy^kz$ for every $k$; these words have distinct lengths, so $L(G)$ is infinite. If the graph has no cycle, a branch of a parse tree, whose consecutive variables are joined by arrows, has no repeated variable, so its height is at most $m$ and the yield has at most $2^{m-1}$ letters: $L(G)$ is finite. So $L(G)$ is infinite if and only if the graph has a cycle, which a depth-first search decides (Tome~C2, Chapter~4).
\end{solution}

\begin{solution}{ex:pdaanbn}
\idea The stack works as a signed counter: it holds only $a$'s or only $b$'s, as many as the surplus of that letter.
\writeup States $q_0$, $q$, $f$ (accepting), stack alphabet $\{a,b,\$\}$. From $q_0$, push $\$$ and go to $q$. In $q$, reading a letter $c\in\{a,b\}$, the PDA may either pop the other letter $\bar c$, if it is on top, or push $c$ (whatever the top): these are the moves $(q,\varepsilon)\in\delta(q,c,\bar c)$ and $(q,c)\in\delta(q,c,\varepsilon)$. Finally $(f,\varepsilon)\in\delta(q,\varepsilon,\$)$.
Let $s$ be the number of $a$'s minus the number of $b$'s on the stack, and $d(u)$ the same count in the input $u$ read so far. Every move reading $a$ increases $s$ by one (it pushes an $a$ or pops a $b$), and every move reading $b$ decreases it by one, so $s=d(u)$ at all times, on every run. To reach $f$ the PDA must pop $\$$, so the stack holds no letter, $s=0$, and the input read is balanced. Conversely, if $d(w)=0$, the run that pops whenever possible keeps on the stack only copies of the letter in surplus, $|d(u)|$ of them; after reading $w$ the stack is $\$$, and the PDA moves to $f$.
\end{solution}

\begin{solution}{ex:linearlength}
\idea A shortest word has a parse tree in which no variable repeats on a branch (otherwise, cut the tree as in Lemma~\ref{lem:cflpumping} to get a shorter word), so its height is at most $m$.
\writeup Let $w\in L(G)$ have a parse tree with the fewest nodes. If some branch carried the same variable twice, replacing the upper subtree by the lower one would give a parse tree with fewer nodes of another word of $L(G)$. So the branches have at most $m$ variables, and $|w|\le2^{m-1}$ by the count in the proof of Lemma~\ref{lem:cflpumping}. The grammar with rules $A_i\to A_{i+1}A_{i+1}$ for $i<m$ and $A_m\to a$, start $A_1$, generates exactly $\{a^{2^{m-1}}\}$, so the bound is attained: grammars can describe very long words with few rules.
\end{solution}

\section*{Sources and further reading}

\begin{itemize}
\item \emph{The stories.} N.~Chomsky, \emph{Three models for the description of language}, IRE Transactions on Information Theory 2 (1956), 113--124; J.~W.~Backus, \emph{The syntax and semantics of the proposed international algebraic language of the Zurich ACM-GAMM conference}, Proceedings of the International Conference on Information Processing, UNESCO, Paris, June 1959, 125--132; P.~Naur (ed.), \emph{Report on the algorithmic language ALGOL~60}, Communications of the ACM 3 (1960), 299--314, whose conditional statements are defined as quoted in the opening; D.~E.~Knuth, \emph{Backus Normal Form vs.\ Backus Naur Form}, Communications of the ACM 7 (1964), 735--736; P.~Z.~Ingerman, \emph{``Pāṇini--Backus form'' suggested}, Communications of the ACM 10 (1967), 137; D.~E.~Knuth, \emph{The remaining trouble spots in ALGOL~60}, Communications of the ACM 10 (1967), 611--618.
\item N.~Chomsky, \emph{On certain formal properties of grammars}, Information and Control 2 (1959), 137--167; Y.~Bar-Hillel, M.~Perles and E.~Shamir, \emph{On formal properties of simple phrase structure grammars}, Zeitschrift für Phonetik, Sprachwissenschaft und Kommunikationsforschung 14 (1961), 143--172; R.~J.~Parikh, \emph{On context-free languages}, Journal of the ACM 13 (1966), 570--581.
\item D.~H.~Younger, \emph{Recognition and parsing of context-free languages in time $n^3$}, Information and Control 10 (1967), 189--208; T.~Kasami, \emph{An efficient recognition and syntax-analysis algorithm for context-free languages}, Air Force Cambridge Research Laboratory report AFCRL-65-758, 1965; J.~Earley, \emph{An efficient context-free parsing algorithm}, Communications of the ACM 13 (1970), 94--102; L.~G.~Valiant, \emph{General context-free recognition in less than cubic time}, Journal of Computer and System Sciences 10 (1975), 308--315; D.~E.~Knuth, \emph{On the translation of languages from left to right}, Information and Control 8 (1965), 607--639.
\item The textbooks of Chapter~\ref{ch:automata}: Hopcroft, Motwani and Ullman for the normal forms and parsing, Sipser for the equivalence of grammars and pushdown automata in the form given here.
\end{itemize}
